\documentclass[10pt,a4paper]{amsart}
\usepackage{amsmath,amssymb,amsthm,graphicx}
% Retained preamble for the Stacks Project chapter morphisms.tex (Morphisms of Schemes).
% The chapter file pinned-original-source.tex carries only \input{preamble} on its first line, so no
% preamble text lives inside the chapter itself.  The Stacks Project uses one single shared preamble for
% every chapter; the snapshot stages it verbatim in sources/stacks-chow/preamble.tex.  The block below is
% copied mechanically, line for line, from that file's lines 14-18 (the Xy-pic setup) and lines 154-217
% (the native theorem declarations together with the mathematical macros).  Nothing was added, removed or
% redefined.  The wrapper class and the amsmath/amssymb/amsthm/graphicx base packages are supplied by the
% assembler, so the class/\externaldocument part of the shared preamble is not repeated here.
% For commutative diagrams we use Xy-pic
\usepackage[all]{xy}

% We use 2cell for 2-commutative diagrams.
\xyoption{2cell}
% Theorem environments.
%
\theoremstyle{plain}
\newtheorem{theorem}[subsection]{Theorem}
\newtheorem{proposition}[subsection]{Proposition}
\newtheorem{lemma}[subsection]{Lemma}

\theoremstyle{definition}
\newtheorem{definition}[subsection]{Definition}
\newtheorem{example}[subsection]{Example}
\newtheorem{exercise}[subsection]{Exercise}
\newtheorem{situation}[subsection]{Situation}

\theoremstyle{remark}
\newtheorem{remark}[subsection]{Remark}
\newtheorem{remarks}[subsection]{Remarks}

\numberwithin{equation}{subsection}

% Macros
%
\def\lim{\mathop{\mathrm{lim}}\nolimits}
\def\colim{\mathop{\mathrm{colim}}\nolimits}
\def\Spec{\mathop{\mathrm{Spec}}}
\def\Hom{\mathop{\mathrm{Hom}}\nolimits}
\def\Ext{\mathop{\mathrm{Ext}}\nolimits}
\def\SheafHom{\mathop{\mathcal{H}\!\mathit{om}}\nolimits}
\def\SheafExt{\mathop{\mathcal{E}\!\mathit{xt}}\nolimits}
\def\Sch{\mathit{Sch}}
\def\Mor{\mathop{\mathrm{Mor}}\nolimits}
\def\Ob{\mathop{\mathrm{Ob}}\nolimits}
\def\Sh{\mathop{\mathit{Sh}}\nolimits}
\def\NL{\mathop{N\!L}\nolimits}
\def\CH{\mathop{\mathrm{CH}}\nolimits}
\def\proetale{{pro\text{-}\acute{e}tale}}
\def\etale{{\acute{e}tale}}
\def\QCoh{\mathit{QCoh}}
\def\Ker{\mathop{\mathrm{Ker}}}
\def\Im{\mathop{\mathrm{Im}}}
\def\Coker{\mathop{\mathrm{Coker}}}
\def\Coim{\mathop{\mathrm{Coim}}}

% Boxtimes
%
\DeclareMathSymbol{\boxtimes}{\mathbin}{AMSa}{"02}

%
% Macros for moduli stacks/spaces
%
\def\QCohstack{\mathcal{QC}\!\mathit{oh}}
\def\Cohstack{\mathcal{C}\!\mathit{oh}}
\def\Spacesstack{\mathcal{S}\!\mathit{paces}}
\def\Quotfunctor{\mathrm{Quot}}
\def\Hilbfunctor{\mathrm{Hilb}}
\def\Curvesstack{\mathcal{C}\!\mathit{urves}}
\def\Polarizedstack{\mathcal{P}\!\mathit{olarized}}
\def\Complexesstack{\mathcal{C}\!\mathit{omplexes}}
% \Pic is the operator that assigns to X its picard group, usage \Pic(X)
% \Picardstack_{X/B} denotes the Picard stack of X over B
% \Picardfunctor_{X/B} denotes the Picard functor of X over B
\def\Pic{\mathop{\mathrm{Pic}}\nolimits}
\def\Picardstack{\mathcal{P}\!\mathit{ic}}
\def\Picardfunctor{\mathrm{Pic}}
\def\Deformationcategory{\mathcal{D}\!\mathit{ef}}

% --- appended 2026-10-12 (Chief): the project own definitions of the reference/slogan/history
% display environments, copied mechanically from the shared Stacks preamble lines 8-12
% (\usepackage{verbatim} plus the three \newenvironment lines, which the chapter preamble keeps
% verbatim from the project preamble). Nothing added, removed or redefined: these environments
% carry no mathematical content and the project itself typesets them as comments. Without these
% three lines a counted statement whose own text uses them fails with Environment slogan undefined.
\usepackage{verbatim}
\newenvironment{reference}{\comment}{\endcomment}
\newenvironment{slogan}{\comment}{\endcomment}
\newenvironment{history}{\comment}{\endcomment}
\usepackage{cleveref}
\crefname{originallemma}{original lemma}{original lemmas}
\Crefname{originallemma}{Original lemma}{Original lemmas}
\makeatletter
\newlabel{morphisms-lemma-rational-map-finite-presentation}{{0BX8}{source}{}{sourcepointer-1}{}}
\newlabel{morphisms-lemma-rational-map-finite-presentation@cref}{{[originallemma][0][]0BX8}{[1][0][]source}{}{}{}}
\makeatother
\crefname{originallemma}{original lemma}{original lemmas}
\Crefname{originallemma}{Original lemma}{Original lemmas}
\makeatletter
\newlabel{lemma-morphism-defined-local-ring}{{0BX6}{source}{}{sourcepointer-2}{}}
\newlabel{lemma-morphism-defined-local-ring@cref}{{[originallemma][0][]0BX6}{[1][0][]source}{}{}{}}
\makeatother
\title{Rational maps to a locally finitely presented scheme are determined by their values at the generic points}
\author{The Stacks Project Authors}
\date{stacks-project:morphisms.tex:local-snapshot-20261012; GNU FDL 1.2 or later}
\begin{document}
\maketitle
\setcounter{subsection}{1}
\noindent\textit{Editorial note.} All delivered mathematical text is the author's own, copied from the pinned Stacks Project chapter file morphisms.tex; no statement or proof step was rewritten, reordered or removed.\par\medskip
\noindent\textit{Editorial note.} The chapter file pinned-original-source.tex carries only \textbackslash{}input\{preamble\} on its first line, so no preamble text lives inside the chapter. The retained file sources/stacks-morphisms/retained-preamble.tex is a mechanical, line-for-line copy of the project's single shared preamble as staged in the snapshot. Nothing was added, removed or redefined, and the excerpt serves only as preamble, never as a span coordinate.\par\medskip
\noindent\textit{Editorial note.} Both counted spans start at the \textbackslash{}begin line of their environment, so the statement and the author's proof are delivered whole; surrounding prose of the section is deliberately outside the counted spans and is carried as standard prerequisites instead.\par\medskip
\noindent\textit{Editorial note.} The citation fields (source\_id, source\_version, authors, source\_url, license.name, source\_title) are identical to those of every other card this lane has produced from this chapter, so the batch selector treats them as one source.\par\medskip
\noindent\textit{Editorial note.} The second paragraph of the proof concludes by citing parts (1) and (2) of an in-chapter lemma, but the reduction to the irreducible case in the first paragraph is the author's own work inside the counted span (separating the components of the source and factoring the correspondence as a product over them), so this is not a citation-only span of the kind refused for two candidates in the sibling half of this chapter.\par\medskip
\noindent\textit{Allowed standard prerequisites (editorial).}\begin{itemize}
\item A morphism is determined by, and locally of finite presentation morphisms are exactly recovered from, the local maps at the generic point. Lemma \textbackslash{}ref\{lemma-morphism-defined-local-ring\}, cited at proof lines 12655-12657. Supplies parts (1), (2) and (3) for an irreducible source.
\item The construction of the map from a rational map by evaluation and pullback of local rings. the definition of a rational map and its local rings at the generic points, used at proof lines 12641-12644. Makes the correspondence of the statement well defined.
\end{itemize}\medskip
\ifdefined\hypertarget\hypertarget{sourcepointer-1}{}\fi
\noindent\textit{Source reference.} Original lemma 0BX8: focused-seeds/AGP/sources/stacks-snapshot-20261011/tags.txt, source lines 14971--14971. The lemma's own canonical tag, quoted so that readers can resolve the label inside the delivered text; it points at the counted statement itself and carries no proof content.\par\medskip
\ifdefined\hypertarget\hypertarget{sourcepointer-2}{}\fi
\noindent\textit{Source reference.} Original lemma 0BX6: focused-seeds/AGP/sources/stacks-snapshot-20261011/tags.txt, source lines 14969--14969. Navigation pointer to a statement the counted proof refers to but does not contain; nothing load-bearing is delegated to it.\par\medskip
\begin{lemma}
\label{lemma-rational-map-finite-presentation}
Let $S$ be a scheme. Let $X$ and $Y$ be schemes over $S$. Assume $X$ has
finitely many irreducible components with generic points
$x_1, \ldots, x_n$. Let $s_i \in S$ be the image of $x_i$.
Consider the map
$$
\left\{
\begin{matrix}
S\text{-rational maps} \\
\text{from }X\text{ to }Y
\end{matrix}
\right\}
\longrightarrow
\left\{
\begin{matrix}
(y_1, \varphi_1, \ldots, y_n, \varphi_n)\text{ where }
y_i \in Y\text{ lies over }s_i\text{ and}\\
\varphi_i : \mathcal{O}_{Y, y_i} \to \mathcal{O}_{X, x_i}
\text{ is a local }\mathcal{O}_{S, s_i}\text{-algebra map}
\end{matrix}
\right\}
$$
which sends $f : U \to Y$ to the $2n$-tuple with
$y_i = f(x_i)$ and $\varphi_i = f^\sharp_{x_i}$. Then
\begin{enumerate}
\item If $Y \to S$ is locally of finite type, then the map is injective.
\item If $Y \to S$ is locally of finite presentation, then the map is bijective.
\item If $Y \to S$ is locally of finite type and $X$ reduced,
then the map is bijective.
\end{enumerate}
\end{lemma}

\begin{proof}
Observe that any dense open of $X$ contains the points $x_i$
so the construction makes sense. To prove (1) or (2)
we may replace $X$ by any dense open. Thus if $Z_1, \ldots, Z_n$
are the irreducible components of $X$, then we may replace
$X$ by $X \setminus \bigcup_{i \not = j} Z_i \cap Z_j$.
After doing this $X$ is the disjoint union of its irreducible
components (viewed as open and closed subschemes). Then both the
right hand side and the left hand side of the arrow are products
over the irreducible components and we reduce to the case where
$X$ is irreducible.

\medskip\noindent
Assume $X$ is irreducible with generic point $x$ lying over $s \in S$.
Part (1) follows from part (1) of
Lemma \ref{lemma-morphism-defined-local-ring}.
Parts (2) and (3) follow from part (2) of the same lemma.
\end{proof}

\end{document}
